% =====================================================================
%  GRADUATION THESIS — Hanoi University of Science and Technology
%  Solar Radio Interferometer Simulation at 610 MHz
%  Author: Nguyen Ba Quan (20224350)
%
%  COMPILE WITH:  pdfLaTeX (Standard in VS Code)
% =====================================================================
% !TeX program = latexmk
\documentclass[12pt,a4paper,oneside]{report}

% ---------- Fonts & encoding (Times New Roman look-alike, full VN) ----
\usepackage[T5]{fontenc}               % Hỗ trợ hiển thị ký tự tiếng Việt
\usepackage[utf8]{inputenc}            % Hỗ trợ gõ mã UTF-8
\usepackage{tgtermes}                  % Times New Roman clone cho pdfLaTeX

% ---------- Page geometry (HUST standard) -----------------------------
\usepackage[a4paper,top=2cm,bottom=2cm,left=3cm,right=2cm]{geometry}

% ---------- Math ------------------------------------------------------
\usepackage{amsmath,amssymb,amsfonts}
\usepackage{bm}
\usepackage{siunitx}
\sisetup{detect-all,per-mode=symbol}
% Force ohm to use math-mode Omega — tgtermes/T5 maps the text-mode glyph to W
\DeclareSIUnit{\ohm}{\ensuremath{\Omega}}

% ---------- Graphics, tables, layout ----------------------------------
\usepackage{graphicx}
\usepackage{overpic}
\graphicspath{{figures/}}
% ---- Draft-safe figures: a missing image shows a placeholder box instead
%      of aborting the build (which would stop VS Code before the ToC pass).
%      Real images are used automatically once present. Remove if undesired.
\makeatletter
\newcommand{\@draftmissingbox}[1]{%
  \fbox{\parbox[c][2.6cm][c]{0.55\linewidth}{\centering\ttfamily
  [missing figure]\\ \detokenize{#1}}}}
\let\@orig@includegraphics\includegraphics
\renewcommand{\includegraphics}[2][]{%
  \begingroup\@tempswafalse
  \IfFileExists{#2}{\@tempswatrue}{}\IfFileExists{#2.pdf}{\@tempswatrue}{}%
  \IfFileExists{#2.png}{\@tempswatrue}{}\IfFileExists{#2.jpg}{\@tempswatrue}{}%
  \IfFileExists{figures/#2}{\@tempswatrue}{}\IfFileExists{figures/#2.pdf}{\@tempswatrue}{}%
  \IfFileExists{figures/#2.png}{\@tempswatrue}{}\IfFileExists{figures/#2.jpg}{\@tempswatrue}{}%
  \if@tempswa\@orig@includegraphics[#1]{#2}\else\@draftmissingbox{#2}\fi
  \endgroup}
\makeatother
\usepackage{float}
\usepackage{booktabs}
\usepackage{array}
\usepackage{tabularx}
\usepackage{multirow}
\usepackage{longtable}
\usepackage{caption}
\usepackage{subcaption}
\captionsetup{font=small,labelfont=bf,labelsep=period}
\usepackage{enumitem}
\usepackage[table]{xcolor}
\usepackage{setspace}
\onehalfspacing
\usepackage{microtype}             % font expansion + character protrusion
\setlength{\emergencystretch}{3em} % last-resort stretch to prevent overfull hbox

% ---------- Headers / footers -----------------------------------------
\usepackage{fancyhdr}
\pagestyle{fancy}
\fancyhf{}
\fancyfoot[C]{\thepage}
\renewcommand{\headrulewidth}{0pt}

% ---------- Section / chapter styling (CHAPTER N. TITLE) ---------------
\usepackage{titlesec}
\titleformat{\chapter}[hang]
  {\normalfont\Large\bfseries\filcenter}
  {CHAPTER \thechapter.}{0.6em}{}
\titlespacing*{\chapter}{0pt}{10pt}{20pt}
\titleformat{\section}{\normalfont\large\bfseries}{\thesection}{0.6em}{}
\titleformat*{\section}{\normalfont\large\bfseries}
\titleformat{\subsection}{\normalfont\normalsize\bfseries}{\thesubsection}{0.6em}{}
\titleformat*{\subsection}{\normalfont\normalsize\bfseries}

% ---------- Hyperlinks -------------------------------------------------
\usepackage[hidelinks]{hyperref}
\hypersetup{
  pdftitle={GRADUATION THESIS},
  pdfauthor={Nguyen Ba Quan},
}

% ---------- Bibliography ----------------------------------------------
\usepackage[numbers,sort&compress]{natbib}

% =====================================================================
%  EDITABLE FIELDS — change these to match your final paperwork
% =====================================================================
\newcommand{\thesistitle}{End-to-End Simulation of a Four-Element UHF Radio Interferometer for Solar Radio Observation at 610~MHz}
\newcommand{\studentname}{Nguyen Ba Quan}
\newcommand{\studentid}{20224350}
\newcommand{\studentemail}{quan.nb224350@sis.hust.edu.vn}
\newcommand{\studentnameVN}{Nguyễn Bá Quân}
\newcommand{\supervisor}{Ha Duyen Trung}
\newcommand{\supervisorVN}{Hà Duyên Trung}
\newcommand{\programme}{Electronics and Telecommunications Engineering (Bachelor of Science)}
\newcommand{\department}{Department of Communication Engineering}
\newcommand{\school}{School of Electrical and Electronic Engineering}
\newcommand{\submitplacedate}{HANOI, 06/2026}

% ---------- Eval-form helpers (eval_forms.tex) ------------------------
\definecolor{LightCyan}{rgb}{0.88,1,1}
\newcolumntype{s}{>{\hsize=0.5\hsize\centering\arraybackslash}X}
\newcolumntype{g}{>{\hsize=3.5\hsize\arraybackslash}X}
\newcolumntype{d}{>{\hsize=0.6\hsize\centering\arraybackslash}X}
\newcommand{\rpt}[2][1]{%
  \ifcase#1\relax
  \or#2%
  \or#2#2%
  \or#2#2#2%
  \or#2#2#2#2%
  \or#2#2#2#2#2%
  \or#2#2#2#2#2#2%
  \fi
}

% Handy math shortcuts used across chapters
\newcommand{\uvw}{(u,v,w)}
\newcommand{\sfu}{\mathrm{sfu}}
\DeclareMathOperator{\sinc}{sinc}
\DeclareMathOperator{\jinc}{jinc}

% Allow \texttt{} to break at underscores so long code identifiers
% (e.g. CLOCK_OFFSET_NS, UFRAME_primary_beam_610MHz) do not overflow the margin.
\let\OrigTexttt\texttt
\renewcommand{\texttt}[1]{%
  \OrigTexttt{\begingroup
    \def\_{\textunderscore\penalty0\relax}%
  #1\endgroup}}

\begin{document}
\renewcommand{\thepage}{\roman{page}}   % front matter in roman numerals
\setcounter{page}{1}

% =====================================================================
%  TITLE PAGE
% =====================================================================
\begin{titlepage}
\thispagestyle{empty}
\begin{center}
{\large\bfseries HANOI UNIVERSITY OF SCIENCE AND TECHNOLOGY}\\[2.2cm]

{\Huge\bfseries GRADUATION THESIS}\\[1.1cm]

{\large\bfseries \thesistitle}\\[1.4cm]

{\large\bfseries \studentname}\\[0.2cm]
{\studentemail}\\[1.0cm]

{\large\bfseries \programme}\\[2.4cm]
\end{center}

\noindent
\begin{tabular}{@{}p{3.0cm}>{\raggedright\arraybackslash}p{7.2cm}p{4.5cm}@{}}
\textbf{Supervisor:} & Assoc. Prof. Dr. Ha Duyen Trung & \\[0.2cm]
                     &             & \rule{4cm}{0.4pt}\\
                     &             & \hfil Signature\hfil\\[0.6cm]
\textbf{Department:} & Department of Communication\newline Engineering & \\[0.3cm]
\textbf{School:}     & School of Electrical and\newline Electronic Engineering & \\
\end{tabular}

\vfill
\begin{center}
{\bfseries \submitplacedate}
\end{center}
\end{titlepage}

% =====================================================================
%  EVALUATION FORMS (HUST standard — Vietnamese)
% =====================================================================
% =====================================================================
%  eval_forms.tex  —  HUST Graduation-Thesis evaluation forms (Vietnamese)
%  Included from main.tex via % =====================================================================
%  eval_forms.tex  —  HUST Graduation-Thesis evaluation forms (Vietnamese)
%  Included from main.tex via % =====================================================================
%  eval_forms.tex  —  HUST Graduation-Thesis evaluation forms (Vietnamese)
%  Included from main.tex via \input{chapters/eval_forms}
%
%  This file is SELF-CONTAINED: it defines its own column types and a
%  local row-helper so it does not depend on macro definitions elsewhere.
%  It compiles with pdfLaTeX + [T5]{fontenc} (Vietnamese) and the packages
%  already loaded in main.tex (tabularx, booktabs/array, multirow,
%  xcolor[table], enumitem).  No siunitx / math is used here.
%
%  Student fields are read from the \student* macros defined in main.tex
%  (\studentnameVN, \studentid, \thesistitle).  Edit those in main.tex.
% =====================================================================

% ---- Local column types (multipliers sum to 4 = number of X columns) -
%  Column 1 (STT)     : narrow, centred
%  Column 2 (Tiêu chí): medium,  left
%  Column 3 (Hướng dẫn): wide,   left, justified
%  Column 4 (Điểm)    : narrow, centred
\newcolumntype{A}{>{\centering\arraybackslash\hsize=0.30\hsize}X}
\newcolumntype{B}{>{\raggedright\arraybackslash\hsize=0.95\hsize}X}
\newcolumntype{G}{>{\raggedright\arraybackslash\hsize=2.30\hsize}X}
\newcolumntype{P}{>{\centering\arraybackslash\hsize=0.45\hsize}X}

% ---- Local helpers ---------------------------------------------------
% Body font for the three forms. If any single form spills one line past
% its page on your machine, change \footnotesize -> \scriptsize here only.
\newcommand{\evalheaderfont}{\footnotesize}
% A grey header row for the four-column criteria tables:
\newcommand{\evalcolhead}{%
  \rowcolor{LightCyan}
  \textbf{STT} & \textbf{Tiêu chí (Điểm tối đa)} &
  \textbf{Hướng dẫn đánh giá tiêu chí} & \textbf{Điểm tiêu chí}\\
}

% =====================================================================
%  These admin forms are front matter: keep them single-spaced and
%  unnumbered, each on its own page, and excluded from the ToC.
% =====================================================================
\begingroup
\singlespacing
\setlength{\extrarowheight}{1pt}\setlength{\tabcolsep}{4pt}
\evalheaderfont

% =====================================================================
%  FORM 1 — DÀNH CHO CÁN BỘ HƯỚNG DẪN
% =====================================================================
\thispagestyle{fancy}
\begin{center}
  {\bfseries ĐẠI HỌC BÁCH KHOA HÀ NỘI}\\[2pt]
  {\bfseries TRƯỜNG ĐIỆN - ĐIỆN TỬ}\\[4pt]
  {\large\bfseries ĐÁNH GIÁ ĐỒ ÁN TỐT NGHIỆP}\\[2pt]
  {\bfseries (DÀNH CHO CÁN BỘ HƯỚNG DẪN)}
\end{center}

\vspace{3pt}
\noindent\textbf{Tên đề tài:} \thesistitle

\vspace{3pt}
\noindent\textbf{Họ tên SV:} \studentnameVN \hfill \textbf{MSSV:} \studentid

\vspace{2pt}
\noindent Cán bộ hướng dẫn 1: \dotfill

\vspace{2pt}
\noindent Cán bộ hướng dẫn 2: \dotfill

\vspace{3pt}
\noindent
\begin{tabularx}{\textwidth}{|A|B|G|P|}
  \hline
  \evalcolhead
  \hline
  \multirow{2}{=}{1}
    & \multirow{2}{=}{\textbf{Thái độ làm việc} (2,5 điểm)}
    & Nghiêm túc, tích cực và chủ động trong quá trình làm ĐATN & \\ \cline{3-4}
    & & Hoàn thành đầy đủ và đúng tiến độ các nội dung được GVHD giao & \\ \hline
  \multirow{2}{=}{2}
    & \multirow{2}{=}{\textbf{Kỹ năng viết quyển ĐATN} (2 điểm)}
    & Trình bày đúng mẫu quy định, bố cục các chương logic và hợp lý:
      bảng biểu, hình ảnh rõ ràng, có tiêu đề, được đánh số thứ tự và được
      giải thích hay đề cập đến trong đồ án, có căn lề, dấu cách sau dấu
      chấm, dấu phẩy, có mở đầu chương và kết luận chương, có liệt kê tài
      liệu tham khảo và có trích dẫn, v.v. & \\ \cline{3-4}
    & & Kỹ năng diễn đạt, phân tích, giải thích, lập luận: cấu trúc câu rõ
        ràng, văn phong khoa học, lập luận logic và có cơ sở, thuật ngữ
        chuyên ngành phù hợp, v.v. & \\ \hline
  \multirow{3}{=}{3}
    & \multirow{3}{=}{\textbf{Nội dung và kết quả đạt được} (5 điểm)}
    & Nêu rõ tính cấp thiết, ý nghĩa khoa học và thực tiễn của đề tài, các
      vấn đề và các giả thuyết, phạm vi ứng dụng của đề tài. Thực hiện đầy
      đủ quy trình nghiên cứu: đặt vấn đề, mục tiêu đề ra, phương pháp
      nghiên cứu/giải quyết vấn đề, kết quả đạt được, đánh giá và kết luận. & \\ \cline{3-4}
    & & Nội dung và kết quả được trình bày một cách logic và hợp lý, được
        phân tích và đánh giá thỏa đáng. Biện luận phân tích kết quả mô
        phỏng/phần mềm/thực nghiệm, so sánh kết quả đạt được với kết quả
        trước đó có liên quan. & \\ \cline{3-4}
    & & Chỉ rõ sự phù hợp giữa kết quả đạt được và mục tiêu ban đầu đề ra,
        đồng thời cung cấp lập luận để đề xuất hướng giải quyết có thể thực
        hiện trong tương lai. Hàm lượng khoa học/độ phức tạp cao, có tính
        mới/tính sáng tạo trong nội dung và kết quả đồ án. & \\ \hline
  \multirow{2}{=}{4}
    & \multirow{2}{=}{\textbf{Điểm thành tích} (1 điểm)}
    & Có bài báo KH được đăng hoặc chấp nhận đăng / đạt giải SV NCKH giải 3
      cấp Trường trở lên / các giải thưởng khoa học trong nước, quốc tế từ
      giải 3 trở lên / có đăng ký bằng phát minh sáng chế. (1 điểm) & \\ \cline{3-4}
    & & Được báo cáo tại hội đồng cấp Trường trong hội nghị SV NCKH nhưng
        không đạt giải từ giải 3 trở lên / đạt giải khuyến khích trong cuộc
        thi khoa học trong nước, quốc tế / kết quả đồ án là sản phẩm ứng
        dụng có tính hoàn thiện cao, yêu cầu khối lượng thực hiện lớn.
        (0,5 điểm) & \\ \hline
  \multicolumn{3}{|l|}{\textbf{Điểm tổng các tiêu chí:}} & \\ \hline
  \multicolumn{3}{|l|}{\textbf{Điểm hướng dẫn:}} & \\ \hline
\end{tabularx}

\vspace{3pt}
\hfill\begin{minipage}{6cm}\centering
  \textbf{Cán bộ hướng dẫn}\\
  (Ký và ghi rõ họ tên)
\end{minipage}

\vspace{2pt}\par\noindent\scriptsize\textit{Điểm từng tiêu chí cho lẻ đến 0,5. Nếu Điểm
tổng các tiêu chí $>$ 10 thì Điểm hướng dẫn làm tròn thành 10.}
\normalsize
\clearpage

% =====================================================================
%  FORM 2 — DÀNH CHO CÁN BỘ THÀNH VIÊN HỘI ĐỒNG
% =====================================================================
\thispagestyle{fancy}
\begin{center}
  {\bfseries ĐẠI HỌC BÁCH KHOA HÀ NỘI}\\[2pt]
  {\bfseries TRƯỜNG ĐIỆN - ĐIỆN TỬ}\\[4pt]
  {\large\bfseries ĐÁNH GIÁ ĐỒ ÁN TỐT NGHIỆP}\\[2pt]
  {\bfseries (DÀNH CHO CÁN BỘ THÀNH VIÊN HỘI ĐỒNG)}
\end{center}

\vspace{3pt}
\noindent Hội đồng số: \dotfill

\vspace{2pt}
\noindent\textbf{Họ tên SV:} \studentnameVN \hfill \textbf{MSSV:} \studentid

\vspace{2pt}
\noindent Cán bộ thành viên HĐ: \dotfill

\vspace{3pt}
\noindent
\begin{tabularx}{\textwidth}{|A|B|G|P|}
  \hline
  \evalcolhead
  \hline
  \multirow{2}{=}{1}
    & \multirow{2}{=}{\textbf{Chất lượng slides/Bản vẽ kỹ thuật} (1,5 điểm)}
    & Sử dụng các minh họa hỗ trợ: hình ảnh, biểu đồ rõ nét và phù hợp,
      dễ hiểu & \\ \cline{3-4}
    & & Không quá nhiều từ, biết sử dụng từ khoá; bố cục logic, có đánh số
        trang & \\ \hline
  \multirow{2}{=}{2}
    & \multirow{2}{=}{\textbf{Kỹ năng thuyết trình} (1,5 điểm)}
    & Tự tin, làm chủ nội dung trình bày, đúng thời gian quy định & \\ \cline{3-4}
    & & Dễ hiểu, dễ theo dõi, lô-gic, lôi cuốn. & \\ \hline
  \multirow{3}{=}{3}
    & \multirow{3}{=}{\textbf{Nội dung và kết quả đạt được} (5 điểm)}
    & Nêu rõ tính cấp thiết, ý nghĩa khoa học và thực tiễn của đề tài, các
      vấn đề và các giả thuyết, phạm vi ứng dụng của đề tài. Thực hiện đầy
      đủ quy trình nghiên cứu: đặt vấn đề, mục tiêu đề ra, phương pháp
      nghiên cứu/giải quyết vấn đề, kết quả đạt được, đánh giá và kết luận. & \\ \cline{3-4}
    & & Nội dung và kết quả được trình bày một cách logic và hợp lý, được
        phân tích và đánh giá thỏa đáng. Biện luận phân tích kết quả mô
        phỏng/phần mềm/thực nghiệm, so sánh kết quả đạt được với kết quả
        trước đó có liên quan. & \\ \cline{3-4}
    & & Chỉ rõ sự phù hợp giữa kết quả đạt được và mục tiêu ban đầu đề ra,
        đồng thời cung cấp lập luận để đề xuất hướng giải quyết có thể thực
        hiện trong tương lai. Hàm lượng khoa học/độ phức tạp cao, có tính
        mới/tính sáng tạo trong nội dung và kết quả đồ án. & \\ \hline
  \multirow{2}{=}{4}
    & \multirow{2}{=}{\textbf{Trả lời câu hỏi} (2,5 điểm)}
    & Trả lời ngắn gọn, chính xác, đi thẳng vào vấn đề của câu hỏi. & \\ \cline{3-4}
    & & Nắm vững kiến thức cơ bản liên quan đến lĩnh vực nghiên cứu/công
        việc của đồ án. & \\ \hline
  \multirow{2}{=}{5}
    & \multirow{2}{=}{\textbf{Điểm thành tích} (1 điểm)}
    & Có bài báo KH được đăng hoặc chấp nhận đăng / đạt giải SV NCKH giải 3
      cấp Trường trở lên / các giải thưởng khoa học trong nước, quốc tế từ
      giải 3 trở lên / có đăng ký bằng phát minh sáng chế. (1 điểm) & \\ \cline{3-4}
    & & Được báo cáo tại hội đồng cấp Trường trong hội nghị SV NCKH nhưng
        không đạt giải từ giải 3 trở lên / đạt giải khuyến khích trong cuộc
        thi khoa học trong nước, quốc tế / kết quả đồ án là sản phẩm ứng
        dụng có tính hoàn thiện cao, yêu cầu khối lượng thực hiện lớn.
        (0,5 điểm) & \\ \hline
  \multicolumn{3}{|l|}{\textbf{Điểm tổng các tiêu chí:}} & \\ \hline
  \multicolumn{3}{|l|}{\textbf{Điểm đánh giá:}} & \\ \hline
\end{tabularx}

\vspace{3pt}
\hfill\begin{minipage}{6cm}\centering
  \textbf{Cán bộ thành viên hội đồng}\\
  (Ký và ghi rõ họ tên)
\end{minipage}

\vspace{2pt}\par\noindent\scriptsize\textit{Điểm từng tiêu chí cho lẻ đến 0,5. Nếu Điểm
tổng các tiêu chí $>$ 10 thì Điểm đánh giá làm tròn thành 10.}
\normalsize
\clearpage

% =====================================================================
%  FORM 3 — DÀNH CHO CÁN BỘ PHẢN BIỆN
% =====================================================================
\thispagestyle{fancy}
\begin{center}
  {\bfseries ĐẠI HỌC BÁCH KHOA HÀ NỘI}\\[2pt]
  {\bfseries TRƯỜNG ĐIỆN - ĐIỆN TỬ}\\[4pt]
  {\large\bfseries ĐÁNH GIÁ ĐỒ ÁN TỐT NGHIỆP}\\[2pt]
  {\bfseries (DÀNH CHO CÁN BỘ PHẢN BIỆN)}
\end{center}

\vspace{3pt}
\noindent\textbf{Tên đề tài:} \thesistitle

\vspace{3pt}
\noindent\textbf{Họ tên SV:} \studentnameVN \hfill \textbf{MSSV:} \studentid

\vspace{2pt}
\noindent Cán bộ phản biện: \dotfill

\vspace{3pt}
\noindent
\begin{tabularx}{\textwidth}{|A|B|G|P|}
  \hline
  \evalcolhead
  \hline
  \multirow{2}{=}{1}
    & \multirow{2}{=}{\textbf{Trình bày quyển ĐATN} (4 điểm)}
    & Đồ án trình bày đúng mẫu quy định, bố cục các chương logic và hợp lý:
      bảng biểu, hình ảnh rõ ràng, có tiêu đề, được đánh số thứ tự và được
      giải thích hay đề cập đến trong đồ án, có căn lề, dấu cách sau dấu
      chấm, dấu phẩy, có mở đầu chương và kết luận chương, có liệt kê tài
      liệu tham khảo và có trích dẫn, v.v. & \\ \cline{3-4}
    & & Kỹ năng diễn đạt, phân tích, giải thích, lập luận: cấu trúc câu rõ
        ràng, văn phong khoa học, lập luận logic và có cơ sở, thuật ngữ
        chuyên ngành phù hợp, v.v. & \\ \hline
  \multirow{3}{=}{2}
    & \multirow{3}{=}{\textbf{Nội dung và kết quả đạt được} (5,5 điểm)}
    & Nêu rõ tính cấp thiết, ý nghĩa khoa học và thực tiễn của đề tài, các
      vấn đề và các giả thuyết, phạm vi ứng dụng của đề tài. Thực hiện đầy
      đủ quy trình nghiên cứu: đặt vấn đề, mục tiêu đề ra, phương pháp
      nghiên cứu/giải quyết vấn đề, kết quả đạt được, đánh giá và kết luận. & \\ \cline{3-4}
    & & Nội dung và kết quả được trình bày một cách logic và hợp lý, được
        phân tích và đánh giá thỏa đáng. Biện luận phân tích kết quả mô
        phỏng/phần mềm/thực nghiệm, so sánh kết quả đạt được với kết quả
        trước đó có liên quan. & \\ \cline{3-4}
    & & Chỉ rõ sự phù hợp giữa kết quả đạt được và mục tiêu ban đầu đề ra,
        đồng thời cung cấp lập luận để đề xuất hướng giải quyết có thể thực
        hiện trong tương lai. Hàm lượng khoa học/độ phức tạp cao, có tính
        mới/tính sáng tạo trong nội dung và kết quả đồ án. & \\ \hline
  \multirow{2}{=}{3}
    & \multirow{2}{=}{\textbf{Điểm thành tích} (1 điểm)}
    & Có bài báo KH được đăng hoặc chấp nhận đăng / đạt giải SV NCKH giải 3
      cấp Trường trở lên / các giải thưởng khoa học trong nước, quốc tế từ
      giải 3 trở lên / có đăng ký bằng phát minh sáng chế. (1 điểm) & \\ \cline{3-4}
    & & Được báo cáo tại hội đồng cấp Trường trong hội nghị SV NCKH nhưng
        không đạt giải từ giải 3 trở lên / đạt giải khuyến khích trong cuộc
        thi khoa học trong nước, quốc tế / kết quả đồ án là sản phẩm ứng
        dụng có tính hoàn thiện cao, yêu cầu khối lượng thực hiện lớn.
        (0,5 điểm) & \\ \hline
  \multicolumn{3}{|l|}{\textbf{Điểm tổng các tiêu chí:}} & \\ \hline
  \multicolumn{3}{|l|}{\textbf{Điểm phản biện:}} & \\ \hline
\end{tabularx}

\vspace{3pt}
\hfill\begin{minipage}{6cm}\centering
  \textbf{Cán bộ phản biện}\\
  (Ký và ghi rõ họ tên)
\end{minipage}

\vspace{2pt}\par\noindent\scriptsize\textit{Điểm từng tiêu chí cho lẻ đến 0,5. Nếu Điểm
tổng các tiêu chí $>$ 10 thì Điểm phản biện làm tròn thành 10.}
\normalsize

\endgroup
\clearpage
%
%  This file is SELF-CONTAINED: it defines its own column types and a
%  local row-helper so it does not depend on macro definitions elsewhere.
%  It compiles with pdfLaTeX + [T5]{fontenc} (Vietnamese) and the packages
%  already loaded in main.tex (tabularx, booktabs/array, multirow,
%  xcolor[table], enumitem).  No siunitx / math is used here.
%
%  Student fields are read from the \student* macros defined in main.tex
%  (\studentnameVN, \studentid, \thesistitle).  Edit those in main.tex.
% =====================================================================

% ---- Local column types (multipliers sum to 4 = number of X columns) -
%  Column 1 (STT)     : narrow, centred
%  Column 2 (Tiêu chí): medium,  left
%  Column 3 (Hướng dẫn): wide,   left, justified
%  Column 4 (Điểm)    : narrow, centred
\newcolumntype{A}{>{\centering\arraybackslash\hsize=0.30\hsize}X}
\newcolumntype{B}{>{\raggedright\arraybackslash\hsize=0.95\hsize}X}
\newcolumntype{G}{>{\raggedright\arraybackslash\hsize=2.30\hsize}X}
\newcolumntype{P}{>{\centering\arraybackslash\hsize=0.45\hsize}X}

% ---- Local helpers ---------------------------------------------------
% Body font for the three forms. If any single form spills one line past
% its page on your machine, change \footnotesize -> \scriptsize here only.
\newcommand{\evalheaderfont}{\footnotesize}
% A grey header row for the four-column criteria tables:
\newcommand{\evalcolhead}{%
  \rowcolor{LightCyan}
  \textbf{STT} & \textbf{Tiêu chí (Điểm tối đa)} &
  \textbf{Hướng dẫn đánh giá tiêu chí} & \textbf{Điểm tiêu chí}\\
}

% =====================================================================
%  These admin forms are front matter: keep them single-spaced and
%  unnumbered, each on its own page, and excluded from the ToC.
% =====================================================================
\begingroup
\singlespacing
\setlength{\extrarowheight}{1pt}\setlength{\tabcolsep}{4pt}
\evalheaderfont

% =====================================================================
%  FORM 1 — DÀNH CHO CÁN BỘ HƯỚNG DẪN
% =====================================================================
\thispagestyle{fancy}
\begin{center}
  {\bfseries ĐẠI HỌC BÁCH KHOA HÀ NỘI}\\[2pt]
  {\bfseries TRƯỜNG ĐIỆN - ĐIỆN TỬ}\\[4pt]
  {\large\bfseries ĐÁNH GIÁ ĐỒ ÁN TỐT NGHIỆP}\\[2pt]
  {\bfseries (DÀNH CHO CÁN BỘ HƯỚNG DẪN)}
\end{center}

\vspace{3pt}
\noindent\textbf{Tên đề tài:} \thesistitle

\vspace{3pt}
\noindent\textbf{Họ tên SV:} \studentnameVN \hfill \textbf{MSSV:} \studentid

\vspace{2pt}
\noindent Cán bộ hướng dẫn 1: \dotfill

\vspace{2pt}
\noindent Cán bộ hướng dẫn 2: \dotfill

\vspace{3pt}
\noindent
\begin{tabularx}{\textwidth}{|A|B|G|P|}
  \hline
  \evalcolhead
  \hline
  \multirow{2}{=}{1}
    & \multirow{2}{=}{\textbf{Thái độ làm việc} (2,5 điểm)}
    & Nghiêm túc, tích cực và chủ động trong quá trình làm ĐATN & \\ \cline{3-4}
    & & Hoàn thành đầy đủ và đúng tiến độ các nội dung được GVHD giao & \\ \hline
  \multirow{2}{=}{2}
    & \multirow{2}{=}{\textbf{Kỹ năng viết quyển ĐATN} (2 điểm)}
    & Trình bày đúng mẫu quy định, bố cục các chương logic và hợp lý:
      bảng biểu, hình ảnh rõ ràng, có tiêu đề, được đánh số thứ tự và được
      giải thích hay đề cập đến trong đồ án, có căn lề, dấu cách sau dấu
      chấm, dấu phẩy, có mở đầu chương và kết luận chương, có liệt kê tài
      liệu tham khảo và có trích dẫn, v.v. & \\ \cline{3-4}
    & & Kỹ năng diễn đạt, phân tích, giải thích, lập luận: cấu trúc câu rõ
        ràng, văn phong khoa học, lập luận logic và có cơ sở, thuật ngữ
        chuyên ngành phù hợp, v.v. & \\ \hline
  \multirow{3}{=}{3}
    & \multirow{3}{=}{\textbf{Nội dung và kết quả đạt được} (5 điểm)}
    & Nêu rõ tính cấp thiết, ý nghĩa khoa học và thực tiễn của đề tài, các
      vấn đề và các giả thuyết, phạm vi ứng dụng của đề tài. Thực hiện đầy
      đủ quy trình nghiên cứu: đặt vấn đề, mục tiêu đề ra, phương pháp
      nghiên cứu/giải quyết vấn đề, kết quả đạt được, đánh giá và kết luận. & \\ \cline{3-4}
    & & Nội dung và kết quả được trình bày một cách logic và hợp lý, được
        phân tích và đánh giá thỏa đáng. Biện luận phân tích kết quả mô
        phỏng/phần mềm/thực nghiệm, so sánh kết quả đạt được với kết quả
        trước đó có liên quan. & \\ \cline{3-4}
    & & Chỉ rõ sự phù hợp giữa kết quả đạt được và mục tiêu ban đầu đề ra,
        đồng thời cung cấp lập luận để đề xuất hướng giải quyết có thể thực
        hiện trong tương lai. Hàm lượng khoa học/độ phức tạp cao, có tính
        mới/tính sáng tạo trong nội dung và kết quả đồ án. & \\ \hline
  \multirow{2}{=}{4}
    & \multirow{2}{=}{\textbf{Điểm thành tích} (1 điểm)}
    & Có bài báo KH được đăng hoặc chấp nhận đăng / đạt giải SV NCKH giải 3
      cấp Trường trở lên / các giải thưởng khoa học trong nước, quốc tế từ
      giải 3 trở lên / có đăng ký bằng phát minh sáng chế. (1 điểm) & \\ \cline{3-4}
    & & Được báo cáo tại hội đồng cấp Trường trong hội nghị SV NCKH nhưng
        không đạt giải từ giải 3 trở lên / đạt giải khuyến khích trong cuộc
        thi khoa học trong nước, quốc tế / kết quả đồ án là sản phẩm ứng
        dụng có tính hoàn thiện cao, yêu cầu khối lượng thực hiện lớn.
        (0,5 điểm) & \\ \hline
  \multicolumn{3}{|l|}{\textbf{Điểm tổng các tiêu chí:}} & \\ \hline
  \multicolumn{3}{|l|}{\textbf{Điểm hướng dẫn:}} & \\ \hline
\end{tabularx}

\vspace{3pt}
\hfill\begin{minipage}{6cm}\centering
  \textbf{Cán bộ hướng dẫn}\\
  (Ký và ghi rõ họ tên)
\end{minipage}

\vspace{2pt}\par\noindent\scriptsize\textit{Điểm từng tiêu chí cho lẻ đến 0,5. Nếu Điểm
tổng các tiêu chí $>$ 10 thì Điểm hướng dẫn làm tròn thành 10.}
\normalsize
\clearpage

% =====================================================================
%  FORM 2 — DÀNH CHO CÁN BỘ THÀNH VIÊN HỘI ĐỒNG
% =====================================================================
\thispagestyle{fancy}
\begin{center}
  {\bfseries ĐẠI HỌC BÁCH KHOA HÀ NỘI}\\[2pt]
  {\bfseries TRƯỜNG ĐIỆN - ĐIỆN TỬ}\\[4pt]
  {\large\bfseries ĐÁNH GIÁ ĐỒ ÁN TỐT NGHIỆP}\\[2pt]
  {\bfseries (DÀNH CHO CÁN BỘ THÀNH VIÊN HỘI ĐỒNG)}
\end{center}

\vspace{3pt}
\noindent Hội đồng số: \dotfill

\vspace{2pt}
\noindent\textbf{Họ tên SV:} \studentnameVN \hfill \textbf{MSSV:} \studentid

\vspace{2pt}
\noindent Cán bộ thành viên HĐ: \dotfill

\vspace{3pt}
\noindent
\begin{tabularx}{\textwidth}{|A|B|G|P|}
  \hline
  \evalcolhead
  \hline
  \multirow{2}{=}{1}
    & \multirow{2}{=}{\textbf{Chất lượng slides/Bản vẽ kỹ thuật} (1,5 điểm)}
    & Sử dụng các minh họa hỗ trợ: hình ảnh, biểu đồ rõ nét và phù hợp,
      dễ hiểu & \\ \cline{3-4}
    & & Không quá nhiều từ, biết sử dụng từ khoá; bố cục logic, có đánh số
        trang & \\ \hline
  \multirow{2}{=}{2}
    & \multirow{2}{=}{\textbf{Kỹ năng thuyết trình} (1,5 điểm)}
    & Tự tin, làm chủ nội dung trình bày, đúng thời gian quy định & \\ \cline{3-4}
    & & Dễ hiểu, dễ theo dõi, lô-gic, lôi cuốn. & \\ \hline
  \multirow{3}{=}{3}
    & \multirow{3}{=}{\textbf{Nội dung và kết quả đạt được} (5 điểm)}
    & Nêu rõ tính cấp thiết, ý nghĩa khoa học và thực tiễn của đề tài, các
      vấn đề và các giả thuyết, phạm vi ứng dụng của đề tài. Thực hiện đầy
      đủ quy trình nghiên cứu: đặt vấn đề, mục tiêu đề ra, phương pháp
      nghiên cứu/giải quyết vấn đề, kết quả đạt được, đánh giá và kết luận. & \\ \cline{3-4}
    & & Nội dung và kết quả được trình bày một cách logic và hợp lý, được
        phân tích và đánh giá thỏa đáng. Biện luận phân tích kết quả mô
        phỏng/phần mềm/thực nghiệm, so sánh kết quả đạt được với kết quả
        trước đó có liên quan. & \\ \cline{3-4}
    & & Chỉ rõ sự phù hợp giữa kết quả đạt được và mục tiêu ban đầu đề ra,
        đồng thời cung cấp lập luận để đề xuất hướng giải quyết có thể thực
        hiện trong tương lai. Hàm lượng khoa học/độ phức tạp cao, có tính
        mới/tính sáng tạo trong nội dung và kết quả đồ án. & \\ \hline
  \multirow{2}{=}{4}
    & \multirow{2}{=}{\textbf{Trả lời câu hỏi} (2,5 điểm)}
    & Trả lời ngắn gọn, chính xác, đi thẳng vào vấn đề của câu hỏi. & \\ \cline{3-4}
    & & Nắm vững kiến thức cơ bản liên quan đến lĩnh vực nghiên cứu/công
        việc của đồ án. & \\ \hline
  \multirow{2}{=}{5}
    & \multirow{2}{=}{\textbf{Điểm thành tích} (1 điểm)}
    & Có bài báo KH được đăng hoặc chấp nhận đăng / đạt giải SV NCKH giải 3
      cấp Trường trở lên / các giải thưởng khoa học trong nước, quốc tế từ
      giải 3 trở lên / có đăng ký bằng phát minh sáng chế. (1 điểm) & \\ \cline{3-4}
    & & Được báo cáo tại hội đồng cấp Trường trong hội nghị SV NCKH nhưng
        không đạt giải từ giải 3 trở lên / đạt giải khuyến khích trong cuộc
        thi khoa học trong nước, quốc tế / kết quả đồ án là sản phẩm ứng
        dụng có tính hoàn thiện cao, yêu cầu khối lượng thực hiện lớn.
        (0,5 điểm) & \\ \hline
  \multicolumn{3}{|l|}{\textbf{Điểm tổng các tiêu chí:}} & \\ \hline
  \multicolumn{3}{|l|}{\textbf{Điểm đánh giá:}} & \\ \hline
\end{tabularx}

\vspace{3pt}
\hfill\begin{minipage}{6cm}\centering
  \textbf{Cán bộ thành viên hội đồng}\\
  (Ký và ghi rõ họ tên)
\end{minipage}

\vspace{2pt}\par\noindent\scriptsize\textit{Điểm từng tiêu chí cho lẻ đến 0,5. Nếu Điểm
tổng các tiêu chí $>$ 10 thì Điểm đánh giá làm tròn thành 10.}
\normalsize
\clearpage

% =====================================================================
%  FORM 3 — DÀNH CHO CÁN BỘ PHẢN BIỆN
% =====================================================================
\thispagestyle{fancy}
\begin{center}
  {\bfseries ĐẠI HỌC BÁCH KHOA HÀ NỘI}\\[2pt]
  {\bfseries TRƯỜNG ĐIỆN - ĐIỆN TỬ}\\[4pt]
  {\large\bfseries ĐÁNH GIÁ ĐỒ ÁN TỐT NGHIỆP}\\[2pt]
  {\bfseries (DÀNH CHO CÁN BỘ PHẢN BIỆN)}
\end{center}

\vspace{3pt}
\noindent\textbf{Tên đề tài:} \thesistitle

\vspace{3pt}
\noindent\textbf{Họ tên SV:} \studentnameVN \hfill \textbf{MSSV:} \studentid

\vspace{2pt}
\noindent Cán bộ phản biện: \dotfill

\vspace{3pt}
\noindent
\begin{tabularx}{\textwidth}{|A|B|G|P|}
  \hline
  \evalcolhead
  \hline
  \multirow{2}{=}{1}
    & \multirow{2}{=}{\textbf{Trình bày quyển ĐATN} (4 điểm)}
    & Đồ án trình bày đúng mẫu quy định, bố cục các chương logic và hợp lý:
      bảng biểu, hình ảnh rõ ràng, có tiêu đề, được đánh số thứ tự và được
      giải thích hay đề cập đến trong đồ án, có căn lề, dấu cách sau dấu
      chấm, dấu phẩy, có mở đầu chương và kết luận chương, có liệt kê tài
      liệu tham khảo và có trích dẫn, v.v. & \\ \cline{3-4}
    & & Kỹ năng diễn đạt, phân tích, giải thích, lập luận: cấu trúc câu rõ
        ràng, văn phong khoa học, lập luận logic và có cơ sở, thuật ngữ
        chuyên ngành phù hợp, v.v. & \\ \hline
  \multirow{3}{=}{2}
    & \multirow{3}{=}{\textbf{Nội dung và kết quả đạt được} (5,5 điểm)}
    & Nêu rõ tính cấp thiết, ý nghĩa khoa học và thực tiễn của đề tài, các
      vấn đề và các giả thuyết, phạm vi ứng dụng của đề tài. Thực hiện đầy
      đủ quy trình nghiên cứu: đặt vấn đề, mục tiêu đề ra, phương pháp
      nghiên cứu/giải quyết vấn đề, kết quả đạt được, đánh giá và kết luận. & \\ \cline{3-4}
    & & Nội dung và kết quả được trình bày một cách logic và hợp lý, được
        phân tích và đánh giá thỏa đáng. Biện luận phân tích kết quả mô
        phỏng/phần mềm/thực nghiệm, so sánh kết quả đạt được với kết quả
        trước đó có liên quan. & \\ \cline{3-4}
    & & Chỉ rõ sự phù hợp giữa kết quả đạt được và mục tiêu ban đầu đề ra,
        đồng thời cung cấp lập luận để đề xuất hướng giải quyết có thể thực
        hiện trong tương lai. Hàm lượng khoa học/độ phức tạp cao, có tính
        mới/tính sáng tạo trong nội dung và kết quả đồ án. & \\ \hline
  \multirow{2}{=}{3}
    & \multirow{2}{=}{\textbf{Điểm thành tích} (1 điểm)}
    & Có bài báo KH được đăng hoặc chấp nhận đăng / đạt giải SV NCKH giải 3
      cấp Trường trở lên / các giải thưởng khoa học trong nước, quốc tế từ
      giải 3 trở lên / có đăng ký bằng phát minh sáng chế. (1 điểm) & \\ \cline{3-4}
    & & Được báo cáo tại hội đồng cấp Trường trong hội nghị SV NCKH nhưng
        không đạt giải từ giải 3 trở lên / đạt giải khuyến khích trong cuộc
        thi khoa học trong nước, quốc tế / kết quả đồ án là sản phẩm ứng
        dụng có tính hoàn thiện cao, yêu cầu khối lượng thực hiện lớn.
        (0,5 điểm) & \\ \hline
  \multicolumn{3}{|l|}{\textbf{Điểm tổng các tiêu chí:}} & \\ \hline
  \multicolumn{3}{|l|}{\textbf{Điểm phản biện:}} & \\ \hline
\end{tabularx}

\vspace{3pt}
\hfill\begin{minipage}{6cm}\centering
  \textbf{Cán bộ phản biện}\\
  (Ký và ghi rõ họ tên)
\end{minipage}

\vspace{2pt}\par\noindent\scriptsize\textit{Điểm từng tiêu chí cho lẻ đến 0,5. Nếu Điểm
tổng các tiêu chí $>$ 10 thì Điểm phản biện làm tròn thành 10.}
\normalsize

\endgroup
\clearpage
%
%  This file is SELF-CONTAINED: it defines its own column types and a
%  local row-helper so it does not depend on macro definitions elsewhere.
%  It compiles with pdfLaTeX + [T5]{fontenc} (Vietnamese) and the packages
%  already loaded in main.tex (tabularx, booktabs/array, multirow,
%  xcolor[table], enumitem).  No siunitx / math is used here.
%
%  Student fields are read from the \student* macros defined in main.tex
%  (\studentnameVN, \studentid, \thesistitle).  Edit those in main.tex.
% =====================================================================

% ---- Local column types (multipliers sum to 4 = number of X columns) -
%  Column 1 (STT)     : narrow, centred
%  Column 2 (Tiêu chí): medium,  left
%  Column 3 (Hướng dẫn): wide,   left, justified
%  Column 4 (Điểm)    : narrow, centred
\newcolumntype{A}{>{\centering\arraybackslash\hsize=0.30\hsize}X}
\newcolumntype{B}{>{\raggedright\arraybackslash\hsize=0.95\hsize}X}
\newcolumntype{G}{>{\raggedright\arraybackslash\hsize=2.30\hsize}X}
\newcolumntype{P}{>{\centering\arraybackslash\hsize=0.45\hsize}X}

% ---- Local helpers ---------------------------------------------------
% Body font for the three forms. If any single form spills one line past
% its page on your machine, change \footnotesize -> \scriptsize here only.
\newcommand{\evalheaderfont}{\footnotesize}
% A grey header row for the four-column criteria tables:
\newcommand{\evalcolhead}{%
  \rowcolor{LightCyan}
  \textbf{STT} & \textbf{Tiêu chí (Điểm tối đa)} &
  \textbf{Hướng dẫn đánh giá tiêu chí} & \textbf{Điểm tiêu chí}\\
}

% =====================================================================
%  These admin forms are front matter: keep them single-spaced and
%  unnumbered, each on its own page, and excluded from the ToC.
% =====================================================================
\begingroup
\singlespacing
\setlength{\extrarowheight}{1pt}\setlength{\tabcolsep}{4pt}
\evalheaderfont

% =====================================================================
%  FORM 1 — DÀNH CHO CÁN BỘ HƯỚNG DẪN
% =====================================================================
\thispagestyle{fancy}
\begin{center}
  {\bfseries ĐẠI HỌC BÁCH KHOA HÀ NỘI}\\[2pt]
  {\bfseries TRƯỜNG ĐIỆN - ĐIỆN TỬ}\\[4pt]
  {\large\bfseries ĐÁNH GIÁ ĐỒ ÁN TỐT NGHIỆP}\\[2pt]
  {\bfseries (DÀNH CHO CÁN BỘ HƯỚNG DẪN)}
\end{center}

\vspace{3pt}
\noindent\textbf{Tên đề tài:} \thesistitle

\vspace{3pt}
\noindent\textbf{Họ tên SV:} \studentnameVN \hfill \textbf{MSSV:} \studentid

\vspace{2pt}
\noindent Cán bộ hướng dẫn 1: \dotfill

\vspace{2pt}
\noindent Cán bộ hướng dẫn 2: \dotfill

\vspace{3pt}
\noindent
\begin{tabularx}{\textwidth}{|A|B|G|P|}
  \hline
  \evalcolhead
  \hline
  \multirow{2}{=}{1}
    & \multirow{2}{=}{\textbf{Thái độ làm việc} (2,5 điểm)}
    & Nghiêm túc, tích cực và chủ động trong quá trình làm ĐATN & \\ \cline{3-4}
    & & Hoàn thành đầy đủ và đúng tiến độ các nội dung được GVHD giao & \\ \hline
  \multirow{2}{=}{2}
    & \multirow{2}{=}{\textbf{Kỹ năng viết quyển ĐATN} (2 điểm)}
    & Trình bày đúng mẫu quy định, bố cục các chương logic và hợp lý:
      bảng biểu, hình ảnh rõ ràng, có tiêu đề, được đánh số thứ tự và được
      giải thích hay đề cập đến trong đồ án, có căn lề, dấu cách sau dấu
      chấm, dấu phẩy, có mở đầu chương và kết luận chương, có liệt kê tài
      liệu tham khảo và có trích dẫn, v.v. & \\ \cline{3-4}
    & & Kỹ năng diễn đạt, phân tích, giải thích, lập luận: cấu trúc câu rõ
        ràng, văn phong khoa học, lập luận logic và có cơ sở, thuật ngữ
        chuyên ngành phù hợp, v.v. & \\ \hline
  \multirow{3}{=}{3}
    & \multirow{3}{=}{\textbf{Nội dung và kết quả đạt được} (5 điểm)}
    & Nêu rõ tính cấp thiết, ý nghĩa khoa học và thực tiễn của đề tài, các
      vấn đề và các giả thuyết, phạm vi ứng dụng của đề tài. Thực hiện đầy
      đủ quy trình nghiên cứu: đặt vấn đề, mục tiêu đề ra, phương pháp
      nghiên cứu/giải quyết vấn đề, kết quả đạt được, đánh giá và kết luận. & \\ \cline{3-4}
    & & Nội dung và kết quả được trình bày một cách logic và hợp lý, được
        phân tích và đánh giá thỏa đáng. Biện luận phân tích kết quả mô
        phỏng/phần mềm/thực nghiệm, so sánh kết quả đạt được với kết quả
        trước đó có liên quan. & \\ \cline{3-4}
    & & Chỉ rõ sự phù hợp giữa kết quả đạt được và mục tiêu ban đầu đề ra,
        đồng thời cung cấp lập luận để đề xuất hướng giải quyết có thể thực
        hiện trong tương lai. Hàm lượng khoa học/độ phức tạp cao, có tính
        mới/tính sáng tạo trong nội dung và kết quả đồ án. & \\ \hline
  \multirow{2}{=}{4}
    & \multirow{2}{=}{\textbf{Điểm thành tích} (1 điểm)}
    & Có bài báo KH được đăng hoặc chấp nhận đăng / đạt giải SV NCKH giải 3
      cấp Trường trở lên / các giải thưởng khoa học trong nước, quốc tế từ
      giải 3 trở lên / có đăng ký bằng phát minh sáng chế. (1 điểm) & \\ \cline{3-4}
    & & Được báo cáo tại hội đồng cấp Trường trong hội nghị SV NCKH nhưng
        không đạt giải từ giải 3 trở lên / đạt giải khuyến khích trong cuộc
        thi khoa học trong nước, quốc tế / kết quả đồ án là sản phẩm ứng
        dụng có tính hoàn thiện cao, yêu cầu khối lượng thực hiện lớn.
        (0,5 điểm) & \\ \hline
  \multicolumn{3}{|l|}{\textbf{Điểm tổng các tiêu chí:}} & \\ \hline
  \multicolumn{3}{|l|}{\textbf{Điểm hướng dẫn:}} & \\ \hline
\end{tabularx}

\vspace{3pt}
\hfill\begin{minipage}{6cm}\centering
  \textbf{Cán bộ hướng dẫn}\\
  (Ký và ghi rõ họ tên)
\end{minipage}

\vspace{2pt}\par\noindent\scriptsize\textit{Điểm từng tiêu chí cho lẻ đến 0,5. Nếu Điểm
tổng các tiêu chí $>$ 10 thì Điểm hướng dẫn làm tròn thành 10.}
\normalsize
\clearpage

% =====================================================================
%  FORM 2 — DÀNH CHO CÁN BỘ THÀNH VIÊN HỘI ĐỒNG
% =====================================================================
\thispagestyle{fancy}
\begin{center}
  {\bfseries ĐẠI HỌC BÁCH KHOA HÀ NỘI}\\[2pt]
  {\bfseries TRƯỜNG ĐIỆN - ĐIỆN TỬ}\\[4pt]
  {\large\bfseries ĐÁNH GIÁ ĐỒ ÁN TỐT NGHIỆP}\\[2pt]
  {\bfseries (DÀNH CHO CÁN BỘ THÀNH VIÊN HỘI ĐỒNG)}
\end{center}

\vspace{3pt}
\noindent Hội đồng số: \dotfill

\vspace{2pt}
\noindent\textbf{Họ tên SV:} \studentnameVN \hfill \textbf{MSSV:} \studentid

\vspace{2pt}
\noindent Cán bộ thành viên HĐ: \dotfill

\vspace{3pt}
\noindent
\begin{tabularx}{\textwidth}{|A|B|G|P|}
  \hline
  \evalcolhead
  \hline
  \multirow{2}{=}{1}
    & \multirow{2}{=}{\textbf{Chất lượng slides/Bản vẽ kỹ thuật} (1,5 điểm)}
    & Sử dụng các minh họa hỗ trợ: hình ảnh, biểu đồ rõ nét và phù hợp,
      dễ hiểu & \\ \cline{3-4}
    & & Không quá nhiều từ, biết sử dụng từ khoá; bố cục logic, có đánh số
        trang & \\ \hline
  \multirow{2}{=}{2}
    & \multirow{2}{=}{\textbf{Kỹ năng thuyết trình} (1,5 điểm)}
    & Tự tin, làm chủ nội dung trình bày, đúng thời gian quy định & \\ \cline{3-4}
    & & Dễ hiểu, dễ theo dõi, lô-gic, lôi cuốn. & \\ \hline
  \multirow{3}{=}{3}
    & \multirow{3}{=}{\textbf{Nội dung và kết quả đạt được} (5 điểm)}
    & Nêu rõ tính cấp thiết, ý nghĩa khoa học và thực tiễn của đề tài, các
      vấn đề và các giả thuyết, phạm vi ứng dụng của đề tài. Thực hiện đầy
      đủ quy trình nghiên cứu: đặt vấn đề, mục tiêu đề ra, phương pháp
      nghiên cứu/giải quyết vấn đề, kết quả đạt được, đánh giá và kết luận. & \\ \cline{3-4}
    & & Nội dung và kết quả được trình bày một cách logic và hợp lý, được
        phân tích và đánh giá thỏa đáng. Biện luận phân tích kết quả mô
        phỏng/phần mềm/thực nghiệm, so sánh kết quả đạt được với kết quả
        trước đó có liên quan. & \\ \cline{3-4}
    & & Chỉ rõ sự phù hợp giữa kết quả đạt được và mục tiêu ban đầu đề ra,
        đồng thời cung cấp lập luận để đề xuất hướng giải quyết có thể thực
        hiện trong tương lai. Hàm lượng khoa học/độ phức tạp cao, có tính
        mới/tính sáng tạo trong nội dung và kết quả đồ án. & \\ \hline
  \multirow{2}{=}{4}
    & \multirow{2}{=}{\textbf{Trả lời câu hỏi} (2,5 điểm)}
    & Trả lời ngắn gọn, chính xác, đi thẳng vào vấn đề của câu hỏi. & \\ \cline{3-4}
    & & Nắm vững kiến thức cơ bản liên quan đến lĩnh vực nghiên cứu/công
        việc của đồ án. & \\ \hline
  \multirow{2}{=}{5}
    & \multirow{2}{=}{\textbf{Điểm thành tích} (1 điểm)}
    & Có bài báo KH được đăng hoặc chấp nhận đăng / đạt giải SV NCKH giải 3
      cấp Trường trở lên / các giải thưởng khoa học trong nước, quốc tế từ
      giải 3 trở lên / có đăng ký bằng phát minh sáng chế. (1 điểm) & \\ \cline{3-4}
    & & Được báo cáo tại hội đồng cấp Trường trong hội nghị SV NCKH nhưng
        không đạt giải từ giải 3 trở lên / đạt giải khuyến khích trong cuộc
        thi khoa học trong nước, quốc tế / kết quả đồ án là sản phẩm ứng
        dụng có tính hoàn thiện cao, yêu cầu khối lượng thực hiện lớn.
        (0,5 điểm) & \\ \hline
  \multicolumn{3}{|l|}{\textbf{Điểm tổng các tiêu chí:}} & \\ \hline
  \multicolumn{3}{|l|}{\textbf{Điểm đánh giá:}} & \\ \hline
\end{tabularx}

\vspace{3pt}
\hfill\begin{minipage}{6cm}\centering
  \textbf{Cán bộ thành viên hội đồng}\\
  (Ký và ghi rõ họ tên)
\end{minipage}

\vspace{2pt}\par\noindent\scriptsize\textit{Điểm từng tiêu chí cho lẻ đến 0,5. Nếu Điểm
tổng các tiêu chí $>$ 10 thì Điểm đánh giá làm tròn thành 10.}
\normalsize
\clearpage

% =====================================================================
%  FORM 3 — DÀNH CHO CÁN BỘ PHẢN BIỆN
% =====================================================================
\thispagestyle{fancy}
\begin{center}
  {\bfseries ĐẠI HỌC BÁCH KHOA HÀ NỘI}\\[2pt]
  {\bfseries TRƯỜNG ĐIỆN - ĐIỆN TỬ}\\[4pt]
  {\large\bfseries ĐÁNH GIÁ ĐỒ ÁN TỐT NGHIỆP}\\[2pt]
  {\bfseries (DÀNH CHO CÁN BỘ PHẢN BIỆN)}
\end{center}

\vspace{3pt}
\noindent\textbf{Tên đề tài:} \thesistitle

\vspace{3pt}
\noindent\textbf{Họ tên SV:} \studentnameVN \hfill \textbf{MSSV:} \studentid

\vspace{2pt}
\noindent Cán bộ phản biện: \dotfill

\vspace{3pt}
\noindent
\begin{tabularx}{\textwidth}{|A|B|G|P|}
  \hline
  \evalcolhead
  \hline
  \multirow{2}{=}{1}
    & \multirow{2}{=}{\textbf{Trình bày quyển ĐATN} (4 điểm)}
    & Đồ án trình bày đúng mẫu quy định, bố cục các chương logic và hợp lý:
      bảng biểu, hình ảnh rõ ràng, có tiêu đề, được đánh số thứ tự và được
      giải thích hay đề cập đến trong đồ án, có căn lề, dấu cách sau dấu
      chấm, dấu phẩy, có mở đầu chương và kết luận chương, có liệt kê tài
      liệu tham khảo và có trích dẫn, v.v. & \\ \cline{3-4}
    & & Kỹ năng diễn đạt, phân tích, giải thích, lập luận: cấu trúc câu rõ
        ràng, văn phong khoa học, lập luận logic và có cơ sở, thuật ngữ
        chuyên ngành phù hợp, v.v. & \\ \hline
  \multirow{3}{=}{2}
    & \multirow{3}{=}{\textbf{Nội dung và kết quả đạt được} (5,5 điểm)}
    & Nêu rõ tính cấp thiết, ý nghĩa khoa học và thực tiễn của đề tài, các
      vấn đề và các giả thuyết, phạm vi ứng dụng của đề tài. Thực hiện đầy
      đủ quy trình nghiên cứu: đặt vấn đề, mục tiêu đề ra, phương pháp
      nghiên cứu/giải quyết vấn đề, kết quả đạt được, đánh giá và kết luận. & \\ \cline{3-4}
    & & Nội dung và kết quả được trình bày một cách logic và hợp lý, được
        phân tích và đánh giá thỏa đáng. Biện luận phân tích kết quả mô
        phỏng/phần mềm/thực nghiệm, so sánh kết quả đạt được với kết quả
        trước đó có liên quan. & \\ \cline{3-4}
    & & Chỉ rõ sự phù hợp giữa kết quả đạt được và mục tiêu ban đầu đề ra,
        đồng thời cung cấp lập luận để đề xuất hướng giải quyết có thể thực
        hiện trong tương lai. Hàm lượng khoa học/độ phức tạp cao, có tính
        mới/tính sáng tạo trong nội dung và kết quả đồ án. & \\ \hline
  \multirow{2}{=}{3}
    & \multirow{2}{=}{\textbf{Điểm thành tích} (1 điểm)}
    & Có bài báo KH được đăng hoặc chấp nhận đăng / đạt giải SV NCKH giải 3
      cấp Trường trở lên / các giải thưởng khoa học trong nước, quốc tế từ
      giải 3 trở lên / có đăng ký bằng phát minh sáng chế. (1 điểm) & \\ \cline{3-4}
    & & Được báo cáo tại hội đồng cấp Trường trong hội nghị SV NCKH nhưng
        không đạt giải từ giải 3 trở lên / đạt giải khuyến khích trong cuộc
        thi khoa học trong nước, quốc tế / kết quả đồ án là sản phẩm ứng
        dụng có tính hoàn thiện cao, yêu cầu khối lượng thực hiện lớn.
        (0,5 điểm) & \\ \hline
  \multicolumn{3}{|l|}{\textbf{Điểm tổng các tiêu chí:}} & \\ \hline
  \multicolumn{3}{|l|}{\textbf{Điểm phản biện:}} & \\ \hline
\end{tabularx}

\vspace{3pt}
\hfill\begin{minipage}{6cm}\centering
  \textbf{Cán bộ phản biện}\\
  (Ký và ghi rõ họ tên)
\end{minipage}

\vspace{2pt}\par\noindent\scriptsize\textit{Điểm từng tiêu chí cho lẻ đến 0,5. Nếu Điểm
tổng các tiêu chí $>$ 10 thì Điểm phản biện làm tròn thành 10.}
\normalsize

\endgroup
\clearpage

% =====================================================================
%  ACKNOWLEDGMENT
% =====================================================================
\chapter*{ACKNOWLEDGMENT}
\addcontentsline{toc}{chapter}{ACKNOWLEDGMENT}
\thispagestyle{fancy}

Completing this graduation thesis marks the end of four years at Hanoi University of Science and Technology. For me, it is a wonderful experience of learning and researching in a creative and competitive environment. 

But to have me today, it is the sacrifice and companionship of my parents, who have always cheered and encouraged me and been by my side, and I am always proud of them. I am grateful and thankful to my parents; no words can express this deep gratitude.

Moreover, I am profoundly grateful to all the teachers who imparted their knowledge and wisdom. I especially want to thank Prof. Han Huy Dung and Prof. Ha Duyen Trung for their patient guidance and the many hours of technical discussions in the SPARC Lab. Prof. Han Huy Dung helped me so much in stepping into the world of radio astronomy, along with Prof. Vo Bich Hien, one of his colleagues in this field. Bringing me into the lab gave me the opportunity to work with the members of the SPARC Laboratory, set up real-time solar radio burst detection antennas, and access the measurement hardware that motivated this simulation work. I consider myself incredibly fortunate to have had their mentorship, guiding me through the challenges and triumphs of my studies. 

Finally, I want to express my commitment to myself: to always stay focused and never give up, even in the most difficult areas of the telecommunication engineering major. Hanoi University of Science and Technology will forever be my second home, the place that helped me grow into a real telecommunication engineer, especially in the field of radio astronomy.


\vspace{1cm}
\hfill\begin{minipage}{6cm}
\centering
\textit{Student}\\
\textit{(Signature and full name)}\\[1.6cm]
\studentname
\end{minipage}

% =====================================================================
%  ABSTRACT
% =====================================================================
\chapter*{ABSTRACT}
\addcontentsline{toc}{chapter}{ABSTRACT}

Space weather driven by solar eruptive activity-flares, coronal mass ejections, and
the geomagnetic storms they cause-can disrupt satellites, GNSS links, and power
grids, as illustrated by the 2022 loss of 38 Starlink satellites to a geomagnetic
storm and by documented degradation of GNSS carrier-to-noise ratios during solar
radio bursts (SRBs). SRBs are an early, ground-observable signature of these
eruptive events, but resolving them with a single large dish is costly, since
angular resolution scales with antenna diameter. Radio interferometry-combining
several small antennas into an aperture-synthesis array, so that resolution instead
depends on the baseline between them-offers a low-cost alternative, a direction
already pursued by projects such as DLITE (the Deployable Low-band Ionosphere and
Transient Experiment), a four-element, low-cost array that has itself detected solar
radio bursts. This thesis pursues the same low-cost, distributed philosophy at
\SI{610}{\mega\hertz}, motivated in part by a preliminary two-element Yagi-Uda field
deployment at Xuan~Mai, Hanoi, that produced only an ambiguous daily trend rather
than a clean interferometric detection, a result traced to phase-incoherent receiver
oscillators, cable insertion loss, local radio-frequency interference, and a noise
floor close to the quiet-Sun signal level. Rather than debug these effects directly
in further hardware iterations, this thesis presents the design and validation of a
complete, physically grounded software simulation of a four-element UHF radio
interferometer, built to separate and quantify such effects in a controlled,
repeatable way, to demonstrate that such an array reproduces the classical
signatures of aperture-synthesis imaging, and to quantify the gap between the
simulation model and the hardware measurement.

The work proceeds in three coupled stages. First, the receiving element, a
commercial UHF Yagi-Uda antenna with a folded-dipole driven element on a
four-bar U-frame reflector, is modelled in 4NEC2 and cross-validated against direct
measurement and an independent finite-element solution. The simulated element exhibits an input impedance of
$Z_\mathrm{in}\approx (108 - j22)\,\si{\ohm}$, a forward gain of about \SI{9.5}{dBi}, a
half-power beamwidth near \SI{52}{\degree}, and a standing-wave ratio of $1.55$ at
\SI{75}{\ohm}; mutual coupling between elements is found to be negligible across the
operational baselines (falling below $-78~\si{dB}$ beyond \SI{5}{\metre}), confirming
that the elements behave as independent sensors. Second, an end-to-end signal-processing pipeline is built in
Python. Real archived solar flux from the SWS Learmonth observatory is converted to
complex antenna voltages through the Friis reception relation with additive
Johnson-Nyquist thermal noise; per-antenna geometric phases are applied using a
standard local-to-equatorial ENU$\rightarrow$XYZ$\rightarrow\uvw$ transformation;
the signals are cross-correlated to form visibilities; and the visibilities are
inverted by a vectorised discrete Fourier transform to a dirty image. Third, the simulated instrument is
exercised over the 20~March~2025 transit at Learmonth station in Australia and over a sweep of
array geometries, including a DLITE-inspired, non-redundant triangle-plus-centre
topology that recovers substantially better North-South resolution than a symmetric
layout using the same four elements. The 20~March~2025 date is also used for a
diagnostic comparison of the simulated solar excess
against a total-power trace from a two-element analogue interferometer deployed
at Hoa~L\d{a}c, Hanoi; a factor of $23.8\times$ gap in SNR is attributed to the
high system temperature of the consumer SDR receiver ($T_{\rm sys,eff}\approx3\,600$~K).

The simulated interferometer reproduces the expected analytic behaviour: a
\mbox{$\sim$14-minute} fringe period for the \SI{8}{\metre} baseline at the spring equinox
($\delta=0^\circ$), a
$\jinc$-shaped visibility decay as a function of baseline length consistent with a
resolved solar disk, and a dirty image equal to the true sky convolved with the
synthesised point-spread function. These results confirm that the simulation faithfully captures the physics from
incident flux to final image at the analytic level.  The hardware comparison further
shows that the primary gap between model and measurement is system temperature, not
signal-chain modelling, a finding that directly guides hardware upgrade priorities.

\paragraph{Keywords:} radio interferometry, aperture synthesis, solar radio bursts,
Yagi-Uda antenna, 4NEC2, van Cittert-Zernike theorem, dirty image, point-spread function.

% =====================================================================
%  TABLE OF CONTENTS / LISTS
% =====================================================================
\renewcommand{\contentsname}{TABLE OF CONTENTS}
\tableofcontents
\clearpage

\renewcommand{\listfigurename}{LIST OF FIGURES}
\listoffigures
\clearpage

\renewcommand{\listtablename}{LIST OF TABLES}
\listoftables
\clearpage

\InputIfFileExists{chapters/abbreviations}{}{\par\noindent\texttt{\detokenize{[chapters/abbreviations.tex not found]}}\par}

% =====================================================================
%  MAIN MATTER
% =====================================================================
\clearpage
\renewcommand{\thepage}{\arabic{page}}
\setcounter{page}{1}

\InputIfFileExists{chapters/intro}{}{\par\noindent\texttt{\detokenize{[chapters/intro.tex not found]}}\par}
\InputIfFileExists{chapters/ch2_theory}{}{\par\noindent\texttt{\detokenize{[chapters/ch2_theory.tex not found]}}\par}
\InputIfFileExists{chapters/ch3_antenna}{}{\par\noindent\texttt{\detokenize{[chapters/ch3_antenna.tex not found]}}\par}
\InputIfFileExists{chapters/ch4_pipeline}{}{\par\noindent\texttt{\detokenize{[chapters/ch4_pipeline.tex not found]}}\par}
\InputIfFileExists{chapters/ch5_results}{}{\par\noindent\texttt{\detokenize{[chapters/ch5_results.tex not found]}}\par}
\InputIfFileExists{chapters/conclusion}{}{\par\noindent\texttt{\detokenize{[chapters/conclusion.tex not found]}}\par}

% =====================================================================
%  BIBLIOGRAPHY
% =====================================================================
\clearpage
\addcontentsline{toc}{chapter}{REFERENCES}
\renewcommand{\bibname}{REFERENCES}
\bibliographystyle{unsrtnat}  % natbib-compatible numbered, unsorted
\bibliography{refs}

% =====================================================================
%  APPENDIX
% =====================================================================
\InputIfFileExists{chapters/appendix_github}{}{\par\noindent\texttt{\detokenize{[chapters/appendix_github.tex not found]}}\par}

\end{document}